\subsection{Actual Taylor-unit gauge: analytic pole classes and their connection}\label{endpoint-analytic_pole--actual-taylor-unit-gauge-analytic-pole-classes-and-their-connection}

13 September 2026. This continues UG1--19 on the original arithmetic packet. The accompanying derivation supplied the exact weighted-residue sequence and the connection on simple-pole lifts. This note proves those constructions in full, calculates their entire finite scaling transport and dual, and compares them with the already proved formal contraction. It does not substitute a new source, a new metric, a simple-root packet, or an assumed arithmetic bound.

\subsubsection{1. Original input and the fixed nonzero-parameter complex}\label{endpoint-analytic_pole--original-input-and-the-fixed-nonzero-parameter-complex}

Keep the monic polynomial ENDPOINTSOURCE3B069B592177DCF6016498D6SLOT0END of degree ENDPOINTSOURCE3B069B592177DCF6016498D6SLOT1END with every selected actual zero order, ENDPOINTSOURCE3B069B592177DCF6016498D6SLOT2END, and the unique polynomial ENDPOINTSOURCE3B069B592177DCF6016498D6SLOT3END of degree less than ENDPOINTSOURCE3B069B592177DCF6016498D6SLOT4END representing the complete Taylor unit ENDPOINTSOURCE3B069B592177DCF6016498D6SLOT5END. Thus ENDPOINTSOURCE3B069B592177DCF6016498D6SLOT6END, and no factor of ENDPOINTSOURCE3B069B592177DCF6016498D6SLOT7END or Taylor coefficient of its unit has been removed. Let ENDPOINTSOURCE3B069B592177DCF6016498D6SLOT8END be the set of distinct zeros of ENDPOINTSOURCE3B069B592177DCF6016498D6SLOT9END, and ENDPOINTSOURCE3B069B592177DCF6016498D6SLOT10END its cardinality; ENDPOINTSOURCE3B069B592177DCF6016498D6SLOT11END when ENDPOINTSOURCE3B069B592177DCF6016498D6SLOT12END is constant. The multiplicities of ENDPOINTSOURCE3B069B592177DCF6016498D6SLOT13END are retained in ENDPOINTSOURCE3B069B592177DCF6016498D6SLOT14END itself throughout the gauge.

Fix the original primitive ENDPOINTSOURCE3B069B592177DCF6016498D6SLOT15END with ENDPOINTSOURCE3B069B592177DCF6016498D6SLOT16END and its specified additive constant. For ENDPOINTSOURCE3B069B592177DCF6016498D6SLOT17END and ENDPOINTSOURCE3B069B592177DCF6016498D6SLOT18END define

ENDPOINTSOURCE3B069B592177DCF6016498D6SLOT19END

Use complexes ENDPOINTSOURCE3B069B592177DCF6016498D6SLOT20END and ENDPOINTSOURCE3B069B592177DCF6016498D6SLOT21END in degrees ENDPOINTSOURCE3B069B592177DCF6016498D6SLOT22END. The equation ENDPOINTSOURCE3B069B592177DCF6016498D6SLOT23END is literal, with the constant in ENDPOINTSOURCE3B069B592177DCF6016498D6SLOT24END retained. Although ENDPOINTSOURCE3B069B592177DCF6016498D6SLOT25END is entire and nowhere zero in ENDPOINTSOURCE3B069B592177DCF6016498D6SLOT26END, it is not asserted rational. The quotient ENDPOINTSOURCE3B069B592177DCF6016498D6SLOT27END consists exactly of finite sums of principal parts at ENDPOINTSOURCE3B069B592177DCF6016498D6SLOT28END: partial fractions gives existence and uniqueness, with the polynomial part discarded only in this explicitly specified quotient.

\subsubsection{2. The entire additional pole cohomology is calculated}\label{endpoint-analytic_pole--the-entire-additional-pole-cohomology-is-calculated}

Define, for ENDPOINTSOURCE3B069B592177DCF6016498D6SLOT29END,

ENDPOINTSOURCE3B069B592177DCF6016498D6SLOT30END

Adding a polynomial changes none of these residues. Also ENDPOINTSOURCE3B069B592177DCF6016498D6SLOT31END by AP1, since a derivative of a meromorphic Laurent series has zero residue.

The exact sequence is

ENDPOINTSOURCE3B069B592177DCF6016498D6SLOT32END

Here is a proof of every assertion. If ENDPOINTSOURCE3B069B592177DCF6016498D6SLOT33END has pole order ENDPOINTSOURCE3B069B592177DCF6016498D6SLOT34END at ENDPOINTSOURCE3B069B592177DCF6016498D6SLOT35END and leading term ENDPOINTSOURCE3B069B592177DCF6016498D6SLOT36END, ENDPOINTSOURCE3B069B592177DCF6016498D6SLOT37END has leading pole term ENDPOINTSOURCE3B069B592177DCF6016498D6SLOT38END. Multiplication by ENDPOINTSOURCE3B069B592177DCF6016498D6SLOT39END has pole order at most ENDPOINTSOURCE3B069B592177DCF6016498D6SLOT40END, so it cannot cancel this term. Thus ENDPOINTSOURCE3B069B592177DCF6016498D6SLOT41END in ENDPOINTSOURCE3B069B592177DCF6016498D6SLOT42END implies that ENDPOINTSOURCE3B069B592177DCF6016498D6SLOT43END has no poles, and ENDPOINTSOURCE3B069B592177DCF6016498D6SLOT44END. This proves injectivity.

For each ENDPOINTSOURCE3B069B592177DCF6016498D6SLOT45END, the rational function ENDPOINTSOURCE3B069B592177DCF6016498D6SLOT46END has residue vector equal to the ENDPOINTSOURCE3B069B592177DCF6016498D6SLOT47END-th coordinate vector. It is analytic at every other point of ENDPOINTSOURCE3B069B592177DCF6016498D6SLOT48END. This proves surjectivity without suppressing the factor ENDPOINTSOURCE3B069B592177DCF6016498D6SLOT49END.

Suppose ENDPOINTSOURCE3B069B592177DCF6016498D6SLOT50END. In a disk at ENDPOINTSOURCE3B069B592177DCF6016498D6SLOT51END containing no other point of ENDPOINTSOURCE3B069B592177DCF6016498D6SLOT52END, the Laurent series of ENDPOINTSOURCE3B069B592177DCF6016498D6SLOT53END has finitely many negative terms and no ENDPOINTSOURCE3B069B592177DCF6016498D6SLOT54END term. Integrating term by term therefore gives a meromorphic primitive ENDPOINTSOURCE3B069B592177DCF6016498D6SLOT55END, with no logarithm. Put ENDPOINTSOURCE3B069B592177DCF6016498D6SLOT56END. Then ENDPOINTSOURCE3B069B592177DCF6016498D6SLOT57END locally. Let ENDPOINTSOURCE3B069B592177DCF6016498D6SLOT58END be the rational function obtained by summing the finite principal parts of the ENDPOINTSOURCE3B069B592177DCF6016498D6SLOT59END at all ENDPOINTSOURCE3B069B592177DCF6016498D6SLOT60END. The difference ENDPOINTSOURCE3B069B592177DCF6016498D6SLOT61END is holomorphic at ENDPOINTSOURCE3B069B592177DCF6016498D6SLOT62END, so ENDPOINTSOURCE3B069B592177DCF6016498D6SLOT63END is holomorphic there. This rational function has no other possible finite poles, hence is a polynomial. Therefore ENDPOINTSOURCE3B069B592177DCF6016498D6SLOT64END in ENDPOINTSOURCE3B069B592177DCF6016498D6SLOT65END, which proves exactness in the middle. Changing any integration constant changes only a holomorphic local function and does not change this construction of the principal part.

All higher pole orders are explicitly observed, rather than ignored:

ENDPOINTSOURCE3B069B592177DCF6016498D6SLOT66END

These formulas give the triangular reduction of each higher pole while retaining every derivative of ENDPOINTSOURCE3B069B592177DCF6016498D6SLOT67END and every derivative of ENDPOINTSOURCE3B069B592177DCF6016498D6SLOT68END that contributes. They explain why there is one quotient coordinate per distinct pole point; they do not replace ENDPOINTSOURCE3B069B592177DCF6016498D6SLOT69END by its radical or discard its nilpotent zero jets.

\subsubsection{3. The original arithmetic classes inject into the analytic extension}\label{endpoint-analytic_pole--the-original-arithmetic-classes-inject-into-the-analytic-extension}

The polynomial complex has ENDPOINTSOURCE3B069B592177DCF6016498D6SLOT70END and ENDPOINTSOURCE3B069B592177DCF6016498D6SLOT71END of dimension ENDPOINTSOURCE3B069B592177DCF6016498D6SLOT72END, with basis ENDPOINTSOURCE3B069B592177DCF6016498D6SLOT73END. Indeed, for a nonzero polynomial of degree ENDPOINTSOURCE3B069B592177DCF6016498D6SLOT74END, the highest term of ENDPOINTSOURCE3B069B592177DCF6016498D6SLOT75END has degree ENDPOINTSOURCE3B069B592177DCF6016498D6SLOT76END, with unchanged leading coefficient. Successive subtraction of such leading terms reduces every polynomial uniquely to degree less than ENDPOINTSOURCE3B069B592177DCF6016498D6SLOT77END. Write this unchanged remainder module as ENDPOINTSOURCE3B069B592177DCF6016498D6SLOT78END.

The same highest-pole argument shows ENDPOINTSOURCE3B069B592177DCF6016498D6SLOT79END of the localized complex is zero: a rational solution of ENDPOINTSOURCE3B069B592177DCF6016498D6SLOT80END has no finite pole, and then the preceding polynomial argument applies. The short exact sequence of the three complexes, together with AP3, gives

ENDPOINTSOURCE3B069B592177DCF6016498D6SLOT81END

One can check injectivity without a long exact sequence. If a polynomial ENDPOINTSOURCE3B069B592177DCF6016498D6SLOT82END equals ENDPOINTSOURCE3B069B592177DCF6016498D6SLOT83END for rational ENDPOINTSOURCE3B069B592177DCF6016498D6SLOT84END, the leading-pole calculation forces ENDPOINTSOURCE3B069B592177DCF6016498D6SLOT85END to be a polynomial. Thus no original polynomial cohomology class becomes a boundary. The kernel and surjectivity assertions follow from AP3 by subtracting ENDPOINTSOURCE3B069B592177DCF6016498D6SLOT86END and retaining the resulting polynomial.

An explicit basis of ENDPOINTSOURCE3B069B592177DCF6016498D6SLOT87END is

ENDPOINTSOURCE3B069B592177DCF6016498D6SLOT88END

For existence, subtract ENDPOINTSOURCE3B069B592177DCF6016498D6SLOT89END. The remainder has zero weighted residues, and AP3 reduces it modulo ENDPOINTSOURCE3B069B592177DCF6016498D6SLOT90END to a polynomial, then to degree less than ENDPOINTSOURCE3B069B592177DCF6016498D6SLOT91END. For independence, applying AP2 to a linear combination forces every pole coefficient to vanish because ENDPOINTSOURCE3B069B592177DCF6016498D6SLOT92END; injectivity in AP5 then gives independence of the polynomial coefficients. Hence ENDPOINTSOURCE3B069B592177DCF6016498D6SLOT93END. The original ENDPOINTSOURCE3B069B592177DCF6016498D6SLOT94END-dimensional arithmetic fibre, with all repeated roots of ENDPOINTSOURCE3B069B592177DCF6016498D6SLOT95END, is present as the explicitly injected summand of this vector-space basis.

\subsubsection{4. Actual parameter connection and both explicit residue frames}\label{endpoint-analytic_pole--actual-parameter-connection-and-both-explicit-residue-frames}

The differential operator ENDPOINTSOURCE3B069B592177DCF6016498D6SLOT96END commutes with ENDPOINTSOURCE3B069B592177DCF6016498D6SLOT97END:

ENDPOINTSOURCE3B069B592177DCF6016498D6SLOT98END

It therefore descends as a connection operator on the analytic family over ENDPOINTSOURCE3B069B592177DCF6016498D6SLOT99END, not as an unsupported multiplication endomorphism on the fixed de Rham quotient. Let ENDPOINTSOURCE3B069B592177DCF6016498D6SLOT100END be the original companion matrix of ENDPOINTSOURCE3B069B592177DCF6016498D6SLOT101END in the polynomial frame, and let ENDPOINTSOURCE3B069B592177DCF6016498D6SLOT102END be its constant-polynomial column. Set

ENDPOINTSOURCE3B069B592177DCF6016498D6SLOT103END

Here each pole basis vector satisfies the exact identity

ENDPOINTSOURCE3B069B592177DCF6016498D6SLOT104END

For the last polynomial vector, reduction of ENDPOINTSOURCE3B069B592177DCF6016498D6SLOT105END uses ENDPOINTSOURCE3B069B592177DCF6016498D6SLOT106END, so the polynomial block is exactly the already specified ENDPOINTSOURCE3B069B592177DCF6016498D6SLOT107END, without a derivative correction. Thus in the full frame AP6 the connection is

ENDPOINTSOURCE3B069B592177DCF6016498D6SLOT108END

All coefficients depend holomorphically on ENDPOINTSOURCE3B069B592177DCF6016498D6SLOT109END for ENDPOINTSOURCE3B069B592177DCF6016498D6SLOT110END. In fact ENDPOINTSOURCE3B069B592177DCF6016498D6SLOT111END itself has no singular ENDPOINTSOURCE3B069B592177DCF6016498D6SLOT112END coefficient; this statement about the matrix does not identify its rank-ENDPOINTSOURCE3B069B592177DCF6016498D6SLOT113END frame with the formal cohomology at ENDPOINTSOURCE3B069B592177DCF6016498D6SLOT114END.

Put ENDPOINTSOURCE3B069B592177DCF6016498D6SLOT115END. The residue map in AP6 is the rectangular matrix

ENDPOINTSOURCE3B069B592177DCF6016498D6SLOT116END

For parameter-dependent representatives direct differentiation proves

ENDPOINTSOURCE3B069B592177DCF6016498D6SLOT117END

Thus the quotient connection in weighted-residue coordinates is ENDPOINTSOURCE3B069B592177DCF6016498D6SLOT118END. The ordinary simple-pole coefficients instead have ENDPOINTSOURCE3B069B592177DCF6016498D6SLOT119END. Their exact transition is multiplication by ENDPOINTSOURCE3B069B592177DCF6016498D6SLOT120END, not deletion of an exponential phase.

For completeness use the second frame ENDPOINTSOURCE3B069B592177DCF6016498D6SLOT121END. Its residue coordinates are the standard coordinate vectors, and

ENDPOINTSOURCE3B069B592177DCF6016498D6SLOT122END

The pole diagonal becomes zero while its coupling to the original polynomial class is ENDPOINTSOURCE3B069B592177DCF6016498D6SLOT123END. This is the same connection under the explicit diagonal frame map ENDPOINTSOURCE3B069B592177DCF6016498D6SLOT124END; no factor has been set to one in the original frame.

\subsubsection{5. Entire finite scaling and its retained extension term}\label{endpoint-analytic_pole--entire-finite-scaling-and-its-retained-extension-term}

Let ENDPOINTSOURCE3B069B592177DCF6016498D6SLOT125END be the original polynomial scaling transport, satisfying

ENDPOINTSOURCE3B069B592177DCF6016498D6SLOT126END

It is the previously proved entire solution, including its full higher terms. Define an integral along the straight segment from ENDPOINTSOURCE3B069B592177DCF6016498D6SLOT127END to ENDPOINTSOURCE3B069B592177DCF6016498D6SLOT128END in the complex ENDPOINTSOURCE3B069B592177DCF6016498D6SLOT129END-plane (equivalently the primitive of its entire integrand):

ENDPOINTSOURCE3B069B592177DCF6016498D6SLOT130END

Differentiating the integral gives ENDPOINTSOURCE3B069B592177DCF6016498D6SLOT131END and ENDPOINTSOURCE3B069B592177DCF6016498D6SLOT132END. Consequently

ENDPOINTSOURCE3B069B592177DCF6016498D6SLOT133END

Existence and uniqueness for this finite polynomial-coefficient matrix ODE prove that these are the actual entire finite connection transports. The block integral retains the coupling into the original polynomial class. Both diagonal blocks are invertible, and the inverse has upper-right block ENDPOINTSOURCE3B069B592177DCF6016498D6SLOT134END. The groupoid law is

ENDPOINTSOURCE3B069B592177DCF6016498D6SLOT135END

It follows either by the shifted ODE and uniqueness, or from the commuting operator's formal flow ENDPOINTSOURCE3B069B592177DCF6016498D6SLOT136END. The latter is coefficientwise in ENDPOINTSOURCE3B069B592177DCF6016498D6SLOT137END when ENDPOINTSOURCE3B069B592177DCF6016498D6SLOT138END is rational; no claim that ENDPOINTSOURCE3B069B592177DCF6016498D6SLOT139END is itself a rational function at a nonzero complex ENDPOINTSOURCE3B069B592177DCF6016498D6SLOT140END is needed. The finite matrix ODE supplies its entire cohomological continuation.

The residue quotient has the exact transport identity

ENDPOINTSOURCE3B069B592177DCF6016498D6SLOT141END

because ENDPOINTSOURCE3B069B592177DCF6016498D6SLOT142END. In weighted residues the entire action is parameter translation; in the unchanged simple-pole frame the factor ENDPOINTSOURCE3B069B592177DCF6016498D6SLOT143END remains. The exponentials and their polynomial extension ENDPOINTSOURCE3B069B592177DCF6016498D6SLOT144END are both displayed, including the case when ENDPOINTSOURCE3B069B592177DCF6016498D6SLOT145END is not real.

\subsubsection{6. Small-loop periods, the actual unit gauge and dual maps}\label{endpoint-analytic_pole--small-loop-periods-the-actual-unit-gauge-and-dual-maps}

For a positively oriented small circle ENDPOINTSOURCE3B069B592177DCF6016498D6SLOT146END enclosing only ENDPOINTSOURCE3B069B592177DCF6016498D6SLOT147END, the exact additional period is

ENDPOINTSOURCE3B069B592177DCF6016498D6SLOT148END

Thus the extra periods are concrete loops, with their orientation and ENDPOINTSOURCE3B069B592177DCF6016498D6SLOT149END factor retained. Every original polynomial image has zero such loop period. For any original rapid-decay contour avoiding ENDPOINTSOURCE3B069B592177DCF6016498D6SLOT150END, the incoming polynomial period is still its old integral, with its original ray orientation and phase.

Choose the finite parts of the original ENDPOINTSOURCE3B069B592177DCF6016498D6SLOT151END rapid-decay contours to avoid every point of ENDPOINTSOURCE3B069B592177DCF6016498D6SLOT152END, including zero when zero belongs to that set. Such detours leave their polynomial integrals unchanged, since those integrands are entire. Different detours for a rational integrand can add the small loops AP19; those are retained rows, not discarded contour ambiguities. At infinity ENDPOINTSOURCE3B069B592177DCF6016498D6SLOT153END has the same decay on the specified rays as in the original polynomial period proof, and the simple-pole functions are ENDPOINTSOURCE3B069B592177DCF6016498D6SLOT154END, so all new integrals converge. On compact parameter subsets in an admitted ENDPOINTSOURCE3B069B592177DCF6016498D6SLOT155END-sector and bounded ENDPOINTSOURCE3B069B592177DCF6016498D6SLOT156END, the negative leading exponential dominates the lower terms and justifies differentiation.

In the unchanged AP6 frame these ENDPOINTSOURCE3B069B592177DCF6016498D6SLOT157END contours followed by the ENDPOINTSOURCE3B069B592177DCF6016498D6SLOT158END small loops give the full period matrix

ENDPOINTSOURCE3B069B592177DCF6016498D6SLOT159END

ENDPOINTSOURCE3B069B592177DCF6016498D6SLOT160END is the original full polynomial period matrix, whose nonzero determinant was proved by the monomial ray matrix and coefficient-connection transport. Hence AP19a is invertible: it realizes all ENDPOINTSOURCE3B069B592177DCF6016498D6SLOT161END analytic classes and all their loop periods. Integrals of ENDPOINTSOURCE3B069B592177DCF6016498D6SLOT162END-exact forms vanish at the decaying endpoints or on closed loops. Differentiating with ENDPOINTSOURCE3B069B592177DCF6016498D6SLOT163END therefore gives

ENDPOINTSOURCE3B069B592177DCF6016498D6SLOT164END

The second identity follows from the same ODE and initial value as AP16, for fixed nonzero ENDPOINTSOURCE3B069B592177DCF6016498D6SLOT165END in the admitted sector. No period matrix is evaluated at ENDPOINTSOURCE3B069B592177DCF6016498D6SLOT166END.

On the same localized complex define the actual gauge differential

ENDPOINTSOURCE3B069B592177DCF6016498D6SLOT167END

The cochain map is multiplication by ENDPOINTSOURCE3B069B592177DCF6016498D6SLOT168END in both degrees. Its inverse is multiplication by ENDPOINTSOURCE3B069B592177DCF6016498D6SLOT169END in ENDPOINTSOURCE3B069B592177DCF6016498D6SLOT170END. The target weight is ENDPOINTSOURCE3B069B592177DCF6016498D6SLOT171END, and the identities

ENDPOINTSOURCE3B069B592177DCF6016498D6SLOT172END

show exact equality of every small-loop and every old rapid-decay period under this map. The target extra class ENDPOINTSOURCE3B069B592177DCF6016498D6SLOT173END need not have a pole as a rational function; its weighted form ENDPOINTSOURCE3B069B592177DCF6016498D6SLOT174END retains the pole. The original image is ENDPOINTSOURCE3B069B592177DCF6016498D6SLOT175END times ENDPOINTSOURCE3B069B592177DCF6016498D6SLOT176END, not the unmodified polynomial frame.

The ordinary algebraic dual of AP5 is the exact sequence

ENDPOINTSOURCE3B069B592177DCF6016498D6SLOT177END

In AP6 the injected column for ENDPOINTSOURCE3B069B592177DCF6016498D6SLOT178END is ENDPOINTSOURCE3B069B592177DCF6016498D6SLOT179END. The dual connection is ENDPOINTSOURCE3B069B592177DCF6016498D6SLOT180END. Its action on that column is ENDPOINTSOURCE3B069B592177DCF6016498D6SLOT181END: the derivative ENDPOINTSOURCE3B069B592177DCF6016498D6SLOT182END cancels ENDPOINTSOURCE3B069B592177DCF6016498D6SLOT183END, and the upper block is zero. Thus the loop-functional injection is explicitly horizontal for the trivial weighted-residue dual. Under the gauge the dual is the inverse transpose of multiplication by ENDPOINTSOURCE3B069B592177DCF6016498D6SLOT184END. This statement concerns the ordinary linear dual; the previous ENDPOINTSOURCE3B069B592177DCF6016498D6SLOT185END-base dagger and moment-line dual retain their separate, already calculated conjugation and weight-line maps. No positive pairing is inferred from AP22, and no original inverse Taylor unit is replaced by ENDPOINTSOURCE3B069B592177DCF6016498D6SLOT186END on a dual.

\subsubsection{7. Exact relation to the formal contraction, including the endpoint}\label{endpoint-analytic_pole--exact-relation-to-the-formal-contraction-including-the-endpoint}

UG uses ENDPOINTSOURCE3B069B592177DCF6016498D6SLOT187END and ENDPOINTSOURCE3B069B592177DCF6016498D6SLOT188END, with coefficientwise completion. On the added pole module ENDPOINTSOURCE3B069B592177DCF6016498D6SLOT189END is invertible because ENDPOINTSOURCE3B069B592177DCF6016498D6SLOT190END. Writing ENDPOINTSOURCE3B069B592177DCF6016498D6SLOT191END for this inverse and ENDPOINTSOURCE3B069B592177DCF6016498D6SLOT192END, its exact formal inverse is the ordered sum ENDPOINTSOURCE3B069B592177DCF6016498D6SLOT193END. That is the contraction proved in UG. AP3 instead fixes ENDPOINTSOURCE3B069B592177DCF6016498D6SLOT194END first and calculates a cokernel of dimension ENDPOINTSOURCE3B069B592177DCF6016498D6SLOT195END, nonzero when ENDPOINTSOURCE3B069B592177DCF6016498D6SLOT196END. For constant ENDPOINTSOURCE3B069B592177DCF6016498D6SLOT197END, ENDPOINTSOURCE3B069B592177DCF6016498D6SLOT198END and the additional pole complex and its residue map are zero; localization adds no analytic or formal classes in that case.

The two operations are related by the original rational complexes and the localization arrow, but the formal completion has no evaluation homomorphism at an arbitrary nonzero complex ENDPOINTSOURCE3B069B592177DCF6016498D6SLOT199END. A formal coefficient sequence may diverge there. Thus the previously displayed formal contraction is not evaluated to claim that AP3's residue classes vanish. AP3 gives exactly the missing cokernel.

There is a concrete endpoint calculation. For ENDPOINTSOURCE3B069B592177DCF6016498D6SLOT200END in ENDPOINTSOURCE3B069B592177DCF6016498D6SLOT201END, ENDPOINTSOURCE3B069B592177DCF6016498D6SLOT202END gives ENDPOINTSOURCE3B069B592177DCF6016498D6SLOT203END. At ENDPOINTSOURCE3B069B592177DCF6016498D6SLOT204END the class of ENDPOINTSOURCE3B069B592177DCF6016498D6SLOT205END modulo ENDPOINTSOURCE3B069B592177DCF6016498D6SLOT206END is the original polynomial class

ENDPOINTSOURCE3B069B592177DCF6016498D6SLOT207END

Indeed multiplying the displayed polynomial by ENDPOINTSOURCE3B069B592177DCF6016498D6SLOT208END gives ENDPOINTSOURCE3B069B592177DCF6016498D6SLOT209END. This is the literal inverse of ENDPOINTSOURCE3B069B592177DCF6016498D6SLOT210END in the original algebra, with every jet of ENDPOINTSOURCE3B069B592177DCF6016498D6SLOT211END retained. For formal ENDPOINTSOURCE3B069B592177DCF6016498D6SLOT212END before the endpoint, ENDPOINTSOURCE3B069B592177DCF6016498D6SLOT213END gives

ENDPOINTSOURCE3B069B592177DCF6016498D6SLOT214END

The successive higher-pole relations AP4 and ordered formal inverse supply the remaining ENDPOINTSOURCE3B069B592177DCF6016498D6SLOT215END coefficients. AP23 is not a rank-(ENDPOINTSOURCE3B069B592177DCF6016498D6SLOT216END) basis at ENDPOINTSOURCE3B069B592177DCF6016498D6SLOT217END. It explicitly describes how these rational lifts are represented in the original rank-ENDPOINTSOURCE3B069B592177DCF6016498D6SLOT218END formal fibre.

The weighted residue factors also retain their parameter behavior: ENDPOINTSOURCE3B069B592177DCF6016498D6SLOT219END. At fixed ENDPOINTSOURCE3B069B592177DCF6016498D6SLOT220END with ENDPOINTSOURCE3B069B592177DCF6016498D6SLOT221END this has an essential singularity at ENDPOINTSOURCE3B069B592177DCF6016498D6SLOT222END. If that numerator vanishes, that particular value is one; its exact derivative is always ENDPOINTSOURCE3B069B592177DCF6016498D6SLOT223END. For ENDPOINTSOURCE3B069B592177DCF6016498D6SLOT224END, ENDPOINTSOURCE3B069B592177DCF6016498D6SLOT225END is independent of ENDPOINTSOURCE3B069B592177DCF6016498D6SLOT226END: it is identically one when ENDPOINTSOURCE3B069B592177DCF6016498D6SLOT227END, and is ENDPOINTSOURCE3B069B592177DCF6016498D6SLOT228END otherwise. Thus no universal singularity or ENDPOINTSOURCE3B069B592177DCF6016498D6SLOT229END-dependence is asserted for every residue coordinate. Even when ENDPOINTSOURCE3B069B592177DCF6016498D6SLOT230END, AP4 on higher pole orders uses ENDPOINTSOURCE3B069B592177DCF6016498D6SLOT231END, whose highest inverse-ENDPOINTSOURCE3B069B592177DCF6016498D6SLOT232END term is ENDPOINTSOURCE3B069B592177DCF6016498D6SLOT233END. Because ENDPOINTSOURCE3B069B592177DCF6016498D6SLOT234END, these contain unbounded inverse-ENDPOINTSOURCE3B069B592177DCF6016498D6SLOT235END powers as the pole order increases near ENDPOINTSOURCE3B069B592177DCF6016498D6SLOT236END. Consequently when ENDPOINTSOURCE3B069B592177DCF6016498D6SLOT237END, AP2 does not define a coefficientwise continuous residue map on UG's completed pole module. For ENDPOINTSOURCE3B069B592177DCF6016498D6SLOT238END its zero map is of course continuous. The obstruction to evaluating the nonzero-pole formal contraction is not replaced by a false assertion that every individual loop factor is singular.

There is also an exact map for the entire displayed endpoint matrix. Let ENDPOINTSOURCE3B069B592177DCF6016498D6SLOT239END have ENDPOINTSOURCE3B069B592177DCF6016498D6SLOT240END-th column the polynomial in AP23, expressed in the original degree-ENDPOINTSOURCE3B069B592177DCF6016498D6SLOT241END frame, and put ENDPOINTSOURCE3B069B592177DCF6016498D6SLOT242END. Direct multiplication of AP23 gives

ENDPOINTSOURCE3B069B592177DCF6016498D6SLOT243END

The map ENDPOINTSOURCE3B069B592177DCF6016498D6SLOT244END carries ENDPOINTSOURCE3B069B592177DCF6016498D6SLOT245END to ENDPOINTSOURCE3B069B592177DCF6016498D6SLOT246END, by the first identity. Thus the endpoint has an explicit invariant ENDPOINTSOURCE3B069B592177DCF6016498D6SLOT247END-dimensional kernel and the original arithmetic quotient, rather than an unexplained rank comparison. Entire matrix exponentiation yields

ENDPOINTSOURCE3B069B592177DCF6016498D6SLOT248END

Here ENDPOINTSOURCE3B069B592177DCF6016498D6SLOT249END is the matrix specialization of AP16; ENDPOINTSOURCE3B069B592177DCF6016498D6SLOT250END then maps it to actual formal-fibre cohomology. Multiplication by the full unit ENDPOINTSOURCE3B069B592177DCF6016498D6SLOT251END after ENDPOINTSOURCE3B069B592177DCF6016498D6SLOT252END gives the gauge endpoint. It commutes with ENDPOINTSOURCE3B069B592177DCF6016498D6SLOT253END since that unit is an element of the original algebra, so the same action diagram holds. The endpoint comparison therefore retains the full unit as well as all jets.

Completed scope: the original arithmetic polynomial cohomology injects into the analytic localization; all additional pole periods, their connection, finite scaling, unit-gauge image and linear dual are calculated. The formal and analytic constructions retain their exact common maps and different coefficient operations. This does not supply a uniform ENDPOINTSOURCE3B069B592177DCF6016498D6SLOT254END-norm bound or an arithmetic Frobenius comparison not constructed in these equations.

\subsubsection{Provenance and verification scope}\label{endpoint-analytic_pole--provenance-and-verification-scope}

Input proof: ACTUAL\_TAYLOR\_UNIT\_FORMAL\_GAUGE.md, SHA256 9ded2486160b76bf65f92d44b73cdd0c59a3ca7e5334f88682e10592d5a45b2f. Original scaling proof after the accepted integer-power-map clarification: ../f1\_scaling\_frobenius/ACTUAL\_TAU\_SCALING\_AND\_EXPONENTIAL\_FLOW.md, SHA256 4e4cb66e3cbbff6379ae389ddaaad65045ac0ed380f17a87ac7c4241c45e6c08. The accompanying derivation supplied AP3's local-primitive construction and AP8--AP13's exact connection in the preceding calculation. This note supplies the full written continuation, AP15--AP18's entire extension transport, AP22's explicit dual and AP23--AP24's endpoint comparison. The independent reviewer supplied the determinant extension AP19a and the endpoint intertwiner AP25--AP26; the text incorporates their full proofs and includes the review's corrected ENDPOINTSOURCE3B069B592177DCF6016498D6SLOT255END residue-factor exception explicitly. No new Lean execution or certified arithmetic-zero computation is claimed.
